\documentclass[11pt]{article}

\usepackage[margin=1in]{geometry}
\usepackage{booktabs}
\usepackage{graphicx}
\usepackage{hyperref}
\usepackage{amsmath}

\title{Two-Factor RASP Budgets for Uneven Decimal Chain Sums\\
\large Alignment $\times$ Carry Clusters, Depth/Head Lower Bounds, and Dual Eval Axes}
\author{YAN Chao}
\date{}

\begin{document}
\maketitle

\begin{abstract}
I factor decimal chain-sum difficulty into two \emph{independent} axes:
\textbf{decimal-point alignment} (fractional unevenness) and \textbf{carry-chain degree}.
Their Cartesian product yields sample clusters $\{A_i\times C_j\}$.
For each cluster I write a specialized RASP program and derive
\emph{minimum depth} and \emph{minimum head} budgets---Transformer-native resource
lower bounds for a faithful align$\to$digitwise-add$\to$carry algorithm.
I then evaluate open models along two complementary trends:
(1)~\textbf{fixed model, vary clusters}; (2)~\textbf{fixed clusters, vary scale}.
I only claim relative difficulty trends and error-type correlations with these budgets,
helping separate architectural constraints from training-distribution effects.
\end{abstract}

\section{Introduction}

Failures on place-value arithmetic are often blamed on missing similar pretraining
examples. RASP~\cite{weiss2021rasp} offers a second account: even ideal data must be
executed under Transformer-native limits on depth (sequential carry) and attention
patterns (alignment selects).

\paragraph{Design.}
I split difficulty into orthogonal factors $A$ (alignment) and $C$ (carry degree),
build independent clusters $A_iC_j$, and for each cluster instantiate a RASP sketch
to read off $\mathrm{depth}_{\min}$ and $\mathrm{heads}_{\min}$.
I then test whether real models track those predictions on two axes---within a model
across clusters, and across scales on fixed clusters.

\paragraph{Scope.}
I only claim relative orders and error-type correlations---not absolute accuracy
forecasts, and not that models have ``learned'' alignment/carry as concepts.

\section{Method}

\subsection{Factors and clusters}
\begin{itemize}
  \item $A_0$: equal fractional lengths (alignment trivial).
  \item $A_1$: uneven fractional lengths (dynamic place binding).
  \item $C_0$: carry-sparse digit bias (short carry runs).
  \item $C_1$: carry-dense digit bias (long carry runs).
\end{itemize}
Clusters: $A_0C_0$, $A_0C_1$, $A_1C_0$, $A_1C_1$ (and extensible grids).

\subsection{Per-cluster RASP $\Rightarrow$ depth/head budgets}
Shared stages: align (if $A_1$), digitwise \texttt{zipmap}, carry scan.
Cluster-specific programs omit or deepen stages; I report
\begin{align*}
\mathrm{depth}_{\min}
&\gtrsim \mathrm{depth}_{\mathrm{align}}(A)
  + 1_{\mathrm{digitwise}}
  + \mathrm{depth}_{\mathrm{carry}}(C,\widehat{\mathrm{run}}), \\
\mathrm{heads}_{\min}
&\gtrsim \#\{\text{distinct select patterns: dot, length, zipmap, pred-carry}\}.
\end{align*}
Predicted easy$\to$hard order typically $A_0C_0 \prec A_1C_0 \prec A_0C_1 \prec A_1C_1$
(exact order follows measured budgets on the generated set).

\subsection{Evaluation axes}
\begin{enumerate}
  \item \textbf{Fixed model, vary clusters:} does accuracy fall / error mass shift as
        $\mathrm{depth}_{\min},\mathrm{heads}_{\min}$ rise?
  \item \textbf{Fixed clusters, vary scale:} do larger models preserve the same relative
        cluster order while lifting absolute accuracy?
\end{enumerate}
Local Ollama pilots ($\approx$1.5--3B) precede larger open/API sweeps on the same JSONL.
Labels: Python exact decimals + Lean $n/10^k$ checks.

\subsection{Metrics}
Exact match; error taxonomy; axis summaries in \texttt{aggregate.json}
(\texttt{fixed\_model\_vary\_clusters}, \texttt{fixed\_cluster\_vary\_models});
join with RASP budgets in \texttt{rasp\_scale\_join.json}.

\section{Results}

\begin{table}[t]
\centering
\caption{Placeholder: cluster accuracies for one fixed model.}
\label{tab:fixed-model}
\begin{tabular}{lcccc}
\toprule
Model & $A_0C_0$ & $A_1C_0$ & $A_0C_1$ & $A_1C_1$ \\
\midrule
Ollama $\approx$3B & -- & -- & -- & -- \\
\bottomrule
\end{tabular}
\end{table}

\begin{table}[t]
\centering
\caption{Placeholder: scale sweep on fixed clusters (relative order).}
\label{tab:fixed-clusters}
\begin{tabular}{lcccc}
\toprule
Scale & $A_0C_0$ & $A_1C_0$ & $A_0C_1$ & $A_1C_1$ \\
\midrule
1.5B / 2B / 3B & -- & -- & -- & -- \\
\bottomrule
\end{tabular}
\end{table}

\section{Discussion}
Agreement with budget-ranked cluster order on both axes supports architecture-linked
difficulty. Absolute gains with scale without order changes suggest capacity lifts the
curve while the Transformer bottleneck shape remains. Order breaks localized to $A$ or
$C$ isolate alignment vs.\ carry-stage failures---useful beyond ``not enough similar data.''

\section{Limitations}
Budgets are faithful-RASP lower bounds, not Tracr-certified circuits. Early pilots use
small local models and smoke-sized clusters.

\section{Conclusion}
By factoring alignment and carry into independent clusters, deriving per-cluster RASP
depth/head budgets, and testing fixed-model vs.\ fixed-cluster trends, I obtain a
diagnostic for whether Transformer constraints predict relative difficulty on uneven
decimal chain sums.

\bibliographystyle{plain}
\bibliography{refs}

\end{document}
